\documentclass[draft]{article}

\usepackage{array}
\usepackage{tabularx}
\usepackage[T1]{fontenc}
\usepackage[utf8x]{inputenc}
\usepackage{graphicx}
\usepackage{hyperref}
\usepackage{booktabs}
\usepackage{xcolor}
\usepackage{xtab}
\usepackage{longtable}
\usepackage{tabu}
\usepackage{colortbl}
\usepackage{varwidth}
\usepackage[a4paper,BCOR12mm,DIV20]{typearea}

\usepackage[ngerman,russian]{babel}
\usepackage{DejaVuSerif}
\newcommand{\ru}[1]{\usefont{T1}{DejaVuSerif-TLF}{m}{n} {\selectlanguage{russian} #1 \selectlanguage{ngerman} \fontfamily{cmr}\selectfont}}
\renewcommand{\rmdefault}{cmr}

\begin{document}

{\fontsize{20pt}{26pt}\selectfont 20 pt} {\fontsize{40pt}{46pt}\selectfont 40 pt} {\fontsize{80pt}{88pt}\selectfont \bfseries 80 pt}

\fbox{
\begin{minipage}{2cm}
minipage
\end{minipage}
}
\fbox{
\begin{varwidth}{2cm}
 varwidth
\end{varwidth}

}

\newcommand{\desc}[2]{ \begin{minipage}{6em} \textbf{#1}\\ \flushright \footnotesize #2 \end{minipage} }

\desc{as}{df}

\begin{longtabu} to \linewidth{lr}
\desc{as}{df} & \desc{gh}{ij}
\end{longtabu}


\begin{tabular}{rlr}
c & \ru{ц}& Москва (\textbf{Moskau}) \\ \midrule
d & \ru{л} & Одеса (Odessa; Ukraine) \\ \midrule
\end{tabular}

{
\newlength{\wRule}
\setlength{\wRule}{1.5pt}
\setlength{\belowrulesep}{0pt}
\setlength{\aboverulesep}{0pt}
\setlength{\cmidrulewidth}{\wRule}
\setlength{\lightrulewidth}{0pt}
\newcommand{\lrule}[1]{\multicolumn{1}{!{\vrule width \wRule}c}{#1}}
\newcommand{\lrrule}[1]{\multicolumn{1}{!{\vrule width \wRule}c!{\vrule width \wRule}}{#1}}

\begin{tabular}{cccc}
 & 14 100 kHz & 18 110 kHz & Dist \\ \cmidrule{2-3}
UN & \lrule{00, 03, …} & \lrrule{+ 10 s}  & 6050 \\
Kanada & \lrule{+ 10 s} & \lrrule{+ 20 s} & 4550 \\
USA & \lrule{+ 20 s} & \lrrule{+ 30 s} & 9050 \\ \cmidrule{2-3}
\end{tabular}
}


{
\newcommand{\freq}[2]{\parbox[t]{6em}{#1\\ \small #2}}

\begin{tabular}{l @{\hspace{4pt}} r @{\hspace{6pt}} p{7cm}}
\toprule
\freq{SLF}{30…300 Hz} & $10^3…10^4$ km & \parbox{7cm}{(Einweg-)Kommunikation zu U-Booten.\\ Elektromagnetische Frequenzen dieser Wellenlänge können \dots} \\ \midrule
\end{tabular}


\begin{tabular}{l @{\hspace{4pt}} r @{\hspace{6pt}} p{7cm}}
\toprule
\freq{SLF}{30…300 Hz} & $10^3…10^4$ km & \parbox[t]{7cm}{(Einweg-)Kommunikation zu U-Booten.\\ Elektromagnetische Frequenzen dieser Wellenlänge können \dots} \\ \midrule
\end{tabular}
}


\newcommand{\wrap}[1]{\begin{minipage}[t]{5.5cm}#1\end{minipage}}
\begin{xtabular}{l p{5cm} p{6cm}}
QRI & Wie ist der Ton meiner Aussendung? & Der Ton Ihrer Aussendung ist 1: gut 2: veränderlich 3: schlecht \\
QRK & Wie ist die Verständlichkeit meiner Zeichen (oder der Zeichen von ...)? & \wrap{Die Verständlichkeit Ihrer Zeichen (oder der Zeichen von ...) ist\\ 1: schlecht\\ 2: mangelhaft\\ 3: ausreichend\\ 4: gut\\ 5: ausgezeichnet} \\
\end{xtabular}


\newlength{\hl}
\settowidth{\hl}{Hello}
Hello is \the\hl\ wide.

Text is \the\textwidth\ wide.

\setlength{\parskip}{\parskip+6pt}

Lines are \the\linewidth\ wide.

Using width of Hello:\\
Hello World.\\
\hspace*{\hl}\ World.

Stretching 1:3: | \hspace{\stretch{1}} 1 \hspace{\stretch{1}} | \hspace{\stretch{3}} 3 \hspace{\stretch{3}} |

Stretching 1:1: | \hfill 1 \hfill | \hfill 1 \hfill |

\newsavebox{\myBox}
\savebox{\myBox}{
Image of the width of this text
}

\newlength{\imgW}
\settowidth{\imgW}{\usebox{\myBox}}

Text width is \the\imgW

\framebox[\textwidth][l]{%
  \includegraphics[width=\imgW]{./png/AvatarGranjow.png}%
  \hspace{2mm}%
  \usebox{\myBox}%
}

\framebox[\textwidth][l]{
 \framebox[.2\textwidth][l]{Hello.}
 \framebox[\dimexpr \textwidth-.2\textwidth][l]{Second part.}
}
\noindent
\setlength{\fboxsep}{0pt}
\framebox[\textwidth][l]{
 \framebox[.2\textwidth][l]{Hello.}
 \framebox[\dimexpr \textwidth-.2\textwidth][l]{Second part.}
}

\noindent
\fbox{
 \framebox[.2\textwidth][l]{Hello.}
 \framebox[\dimexpr \textwidth-.2\textwidth][l]{Second part.}
}

\noindent
\begin{minipage}{\textwidth}
 \parbox{.2\textwidth}{Hello.}
 \parbox{\dimexpr \textwidth-.2\textwidth}{Second part.}
\end{minipage}


\noindent
\begin{minipage}{\textwidth}
 \parbox[t]{1.2cm}{ \raisebox{-3mm}{\includegraphics[height=6mm]{./png/Amfu-Icon-Buch.png}} }
 \parbox[t]{\dimexpr \textwidth-1.2cm}{PDF-Dokument über Typografie: «typokurz -- Einige wichtige typografische Regeln» von Christoph Bier, \href{http://www.zvisionwelt.de}{zvisionwelt.de} }
\end{minipage}



\begin{minipage}[t]{\textwidth}
\colorbox{red}{
\begin{minipage}[t]{.45\textwidth}
\begin{tabular}{r @{---} l l}
1.810 & 2.000 & MHz \\ \midrule
3.500 & 3.800 & MHz \\ \midrule
\end{tabular}
\vspace{\stretch{1}}
\end{minipage}
}
\colorbox{blue}{
\begin{minipage}[t]{.45\textwidth}
\begin{tabular}{r @{---} l l}
 135.7 & 137.8 & kHz \\ \midrule
1.810 & 2.000 & MHz \\ \midrule
3.500 & 3.800 & MHz \\ \midrule
7.000 & 7.200 & MHz \\ \midrule
10.100 & 10.150 & MHz \\ \midrule
14.000 & 14.350 & MHz \\ \midrule
18.068 & 18.168 & MHz \\ \midrule
\end{tabular}
\end{minipage}
}
\end{minipage}



\end{document}
